% Abstract and index terms (final report #63).
% Draft: the speed-ups are the cycle model's (docs/spec.md section 11). At assembly,
% replace them with the measured values (#69, #64) and add the synthesis results (#75).

\begin{abstract}
Matrix multiplication dominates the run time of neural-network inference and many signal-processing workloads, and on general-purpose processors its cost is set by data movement rather than arithmetic. This paper presents the final design and evaluation of a matrix-multiplication accelerator built on an output-stationary systolic array. An $N \times N$ grid of processing elements, $4 \times 4$ by default, multiplies signed 8-bit fixed-point operands and accumulates each element of the result in place in a 32-bit register. The host sets the inner length of each run, from 1 to 16, and can start a run that adds onto the previous result instead of clearing it. These two additions leave the array itself unchanged, yet they let a host multiply matrices of any size, one $N \times N$ tile of the result and one chunk of the inner dimension at a time. The design is synthesizable, parameterized Verilog-2005 with a written timing contract: a run with inner length $k$ is busy for $k + 2N$ cycles. To measure the benefit of the array rather than derive it, we built a sequential baseline with one multiply-accumulate unit and the same host interface, which is busy for $N^2 k$ cycles. The array needs 5.3 times fewer busy cycles than the baseline for $k = 4$ and 10.7 times fewer for $k = 16$, but loading the operands through one-element ports limits the end-to-end speed-up of an $8 \times 8 \times 8$ product to 2.7. Both designs are synthesized for a Lattice ECP5 FPGA with the open-source Yosys and nextpnr tools, and we report their resources and maximum clock frequency. Self-checking testbenches for every module, a Python golden model and continuous integration verify the design, including inner lengths other than $N$ and tiled products larger than the array.
\end{abstract}

\begin{IEEEkeywords}
systolic array, matrix multiplication, hardware accelerator, output-stationary dataflow, tiling, fixed-point arithmetic, FPGA synthesis, Verilog
\end{IEEEkeywords}
